% GENERATED FILE. Edit canonical lessons, then run npm run generate:books.
% canonical-source-hash: fnv1a64:7633b1b01410d514
% canonical-lessons: MR-C474-nitnetka, MR-C474-gabala, MR-C474-vendhla, MR-C474-pasara, MR-C474-sista

\chapter{Tidy and Scatterbrained}
\label{ch:mr-tidy-and-scatterbrained}
\hlchaptermodality{474}

\begin{chapteropening}
\textbf{By the end of this chapter:} \emph{I can describe a tidy, sloppy or scatterbrained person, a mess and discipline.}

Complete the last lesson of chapter 474: I can describe a tidy, sloppy or scatterbrained person, a mess and discipline.
\end{chapteropening}

\section[\texorpdfstring{\mr{नीटनेटका} (nīṭneṭkā)}{नीटनेटका (nīṭneṭkā)}]{\mr{नीटनेटका} (nīṭneṭkā) — tidy; neat}

\label{lesson:MR-C474-nitnetka}



\begin{quote}
Before the new one: say the Marathi for amusement, then the Marathi for to tease.
\end{quote}

\subsection*{You'll want to know: \mr{नीटनेटका}}
\textbf{\mr{नीटनेटका}} — \emph{nīṭneṭkā} — \textquotedblleft{}tidy; neat\textquotedblright{}.

\begin{quote}
\mr{तिची} \mr{खोली} \mr{नेहमी} \mr{नीटनेटकी} \mr{असते}.
\end{quote}

\begin{quote}
\emph{ticī kholī nehmī nīṭneṭkī aste.}
\end{quote}

\begin{quote}
Her room is always neat and tidy.
\end{quote}

Agrees: \mr{नीटनेटका}, \mr{नीटनेटकी}, \mr{नीटनेटकं}.

\begin{cousinweb}
\textbf{\mr{नीटनेटका}} (\emph{nīṭneṭkā}) is a Marathi pair-word, \mr{नीट}, \textquotedblleft{}properly\textquotedblright{}, with \mr{नेटका}, \textquotedblleft{}trim\textquotedblright{}; the origin of each part is not settled.
\end{cousinweb}

\begin{grammarlens}[title={What you've built}]
One more word for describing what you lived and felt.
\end{grammarlens}

\subsection*{Your turn}
\begin{itemize}
\raggedright
  \item \emph{Say it:} \emph{nīṭneṭkā}
  \item \emph{Say it:} \emph{nīṭneṭkā}, once more
  \item \emph{Say it:} \emph{ticī kholī nehmī nīṭneṭkī aste.} — Her room is always neat and tidy.
  \item \emph{Recall:} say the Marathi for amusement, then the Marathi for to tease, then say \emph{nīṭneṭkā} again
\end{itemize}

\begin{tcolorbox}[breakable,colback=teal!4,colframe=teal!35!black,title={Before you move on}]
What does \mr{नीटनेटका} mean? (\textquotedblleft{}Tidy; neat\textquotedblright{}.) Say it once more.
\end{tcolorbox}
\section[\texorpdfstring{\mr{गबाळा} (gabāḷā)}{गबाळा (gabāḷā)}]{\mr{गबाळा} (gabāḷā) — sloppy; untidy}

\label{lesson:MR-C474-gabala}



\begin{quote}
Before the new one: say the Marathi for to tease, then the Marathi for tidy.
\end{quote}

\subsection*{You'll want to know: \mr{गबाळा}}
\textbf{\mr{गबाळा}} — \emph{gabāḷā} — \textquotedblleft{}sloppy; untidy\textquotedblright{}.

\begin{quote}
\mr{गबाळे} \mr{कपडे} \mr{घालून} \mr{तो} \mr{मुलाखतीला} \mr{गेला}.
\end{quote}

\begin{quote}
\emph{gabāḷe kapḍe ghālūn to mulākhtīlā gelā.}
\end{quote}

\begin{quote}
He went to the interview in sloppy clothes.
\end{quote}

Agrees; the opposite of \mr{नीटनेटका}.

\begin{cousinweb}
\textbf{\mr{गबाळा}} (\emph{gabāḷā}) is a Marathi adjective whose origin is not settled.
\end{cousinweb}

\begin{grammarlens}[title={What you've built}]
One more word for describing what you lived and felt.
\end{grammarlens}

\subsection*{Your turn}
\begin{itemize}
\raggedright
  \item \emph{Say it:} \emph{gabāḷā}
  \item \emph{Say it:} \emph{gabāḷā}, once more
  \item \emph{Say it:} \emph{gabāḷe kapḍe ghālūn to mulākhtīlā gelā.} — He went to the interview in sloppy clothes.
  \item \emph{Recall:} say the Marathi for to tease, then the Marathi for tidy, then say \emph{gabāḷā} again
\end{itemize}

\begin{tcolorbox}[breakable,colback=teal!4,colframe=teal!35!black,title={Before you move on}]
What does \mr{गबाळा} mean? (\textquotedblleft{}Sloppy; untidy\textquotedblright{}.) Say it once more.
\end{tcolorbox}
\section[\texorpdfstring{\mr{वेंधळा} (vendhḷā)}{वेंधळा (vendhḷā)}]{\mr{वेंधळा} (vendhḷā) — scatterbrained}

\label{lesson:MR-C474-vendhla}



\begin{quote}
Before the new one: say the Marathi for tidy, then the Marathi for sloppy.
\end{quote}

\subsection*{You'll want to know: \mr{वेंधळा}}
\textbf{\mr{वेंधळा}} — \emph{vendhḷā} — \textquotedblleft{}scatterbrained\textquotedblright{}.

\begin{quote}
\mr{वेंधळा} \mr{अरुण} \mr{पुन्हा} \mr{किल्ली} \mr{विसरला}.
\end{quote}

\begin{quote}
\emph{vendhḷā aruṇ punhā killī visarlā.}
\end{quote}

\begin{quote}
Scatterbrained Arun forgot the key again.
\end{quote}

Agrees; \mr{वेंधळेपणा}, scatterbrained ways.

\begin{cousinweb}
\textbf{\mr{वेंधळा}} (\emph{vendhḷā}) is a Marathi adjective whose origin is not settled.
\end{cousinweb}

\begin{grammarlens}[title={What you've built}]
One more word for describing what you lived and felt.
\end{grammarlens}

\subsection*{Your turn}
\begin{itemize}
\raggedright
  \item \emph{Say it:} \emph{vendhḷā}
  \item \emph{Say it:} \emph{vendhḷā}, once more
  \item \emph{Say it:} \emph{vendhḷā aruṇ punhā killī visarlā.} — Scatterbrained Arun forgot the key again.
  \item \emph{Recall:} say the Marathi for tidy, then the Marathi for sloppy, then say \emph{vendhḷā} again
\end{itemize}

\begin{tcolorbox}[breakable,colback=teal!4,colframe=teal!35!black,title={Before you move on}]
What does \mr{वेंधळा} mean? (\textquotedblleft{}Scatterbrained\textquotedblright{}.) Say it once more.
\end{tcolorbox}
\section[\texorpdfstring{\mr{पसारा} (pasārā)}{पसारा (pasārā)}]{\mr{पसारा} (pasārā) — a mess; clutter}

\label{lesson:MR-C474-pasara}



\begin{quote}
Before the new one: say the Marathi for sloppy, then the Marathi for scatterbrained.
\end{quote}

\subsection*{You'll want to know: \mr{पसारा}}
\textbf{\mr{पसारा}} — \emph{pasārā} — \textquotedblleft{}a mess; clutter\textquotedblright{}.

\begin{quote}
\mr{मुलांच्या} \mr{खोलीत} \mr{खूप} \mr{पसारा} \mr{आहे}.
\end{quote}

\begin{quote}
\emph{mulāñcyā kholīt khūp pasārā āhe.}
\end{quote}

\begin{quote}
There is a big mess in the children's room.
\end{quote}

Masculine.

Masculine: \mr{पसारा} \mr{करणे}, to make a mess; \mr{पसारा} \mr{आवरणे}, to clear it up.

\begin{cousinweb}
\textbf{\mr{पसारा}} (\emph{pasārā}) is a tadbhava: CDIAL lists Marathi \emph{pasārā}, \textquotedblleft{}stretching out, expansion\textquotedblright{}, under Sanskrit \emph{prasāra}, \textquotedblleft{}extension\textquotedblright{} (CDIAL 8835). Things spread out all over a room make a mess.
\end{cousinweb}

\begin{grammarlens}[title={What you've built}]
One more word for describing what you lived and felt.
\end{grammarlens}

\subsection*{Your turn}
\begin{itemize}
\raggedright
  \item \emph{Say it:} \emph{pasārā}
  \item \emph{Say it:} \emph{pasārā}, once more
  \item \emph{Say it:} \emph{mulāñcyā kholīt khūp pasārā āhe.} — There is a big mess in the children's room.
  \item \emph{Recall:} say the Marathi for sloppy, then the Marathi for scatterbrained, then say \emph{pasārā} again
\end{itemize}

\begin{tcolorbox}[breakable,colback=teal!4,colframe=teal!35!black,title={Before you move on}]
What does \mr{पसारा} mean? (\textquotedblleft{}A mess; clutter\textquotedblright{}.) Say it once more.
\end{tcolorbox}
\section[\texorpdfstring{\mr{शिस्त} (śista)}{शिस्त (śista)}]{\mr{शिस्त} (śista) — discipline}

\label{lesson:MR-C474-sista}



\begin{quote}
Before the new one: say the Marathi for scatterbrained, then the Marathi for a mess.
\end{quote}

\subsection*{You'll want to know: \mr{शिस्त}}
\textbf{\mr{शिस्त}} — \emph{śista} — \textquotedblleft{}discipline\textquotedblright{}.

\begin{quote}
\mr{सैन्यात} \mr{कडक} \mr{शिस्त} \mr{असते}.
\end{quote}

\begin{quote}
\emph{sainyāt kaḍak śista aste.}
\end{quote}

\begin{quote}
There is strict discipline in the army.
\end{quote}

Feminine.

\mr{शिस्त} \mr{लावणे}, to instil discipline; \mr{शिस्तीने}, in an orderly way.

\begin{cousinweb}
\textbf{\mr{शिस्त}} (\emph{śista}) is a Marathi noun whose origin is not settled.
\end{cousinweb}

\begin{grammarlens}[title={What you've built}]
One more word for describing what you lived and felt.
\end{grammarlens}

\subsection*{Your turn}
\begin{itemize}
\raggedright
  \item \emph{Say it:} \emph{śista}
  \item \emph{Say it:} \emph{śista}, once more
  \item \emph{Say it:} \emph{sainyāt kaḍak śista aste.} — There is strict discipline in the army.
  \item \emph{Recall:} say the Marathi for scatterbrained, then the Marathi for a mess, then say \emph{śista} again
\end{itemize}

\begin{tcolorbox}[breakable,colback=teal!4,colframe=teal!35!black,title={Before you move on}]
What does \mr{शिस्त} mean? (\textquotedblleft{}Discipline\textquotedblright{}.) Say it once more.
\end{tcolorbox}
